\chapter{绪论}\label{ch:intro}

\section{研究背景与意义}
推荐系统是现代信息服务的核心基础设施：电商、内容平台与社交网络
均依赖个性化排序从海量候选中筛选用户可能感兴趣的内容。随着深度
学习的发展，序列推荐模型——从基于循环神经网络的
GRU4Rec~\cite{hidasi2016gru4rec}、基于自注意力的
SASRec~\cite{kang2018sasrec} 与 BERT4Rec~\cite{sun2019bert4rec}，
到图协同模型 NGCF~\cite{wang2019ngcf} 与对比学习方法
SimGCL~\cite{yu2022simgcl}——已能从用户历史交互中学习细粒度的
偏好表示。然而，监督信号的获取始终是推荐系统的结构性瓶颈：真实
交互数据稀疏、新物品/新用户冷启动严重、人工标注隐式反馈的代价
高昂。

与此同时，大语言模型领域出现了一类新的训练范式：\textbf{递归自
改进（Recursive Self-Improvement, RSI）}——模型\emph{自我生成}训
练数据、\emph{自我评估}数据质量、\emph{自我迭代}优化参数，形成
不依赖外部标注的闭环。代表性工作包括：Self-Refine~\cite%
{madaan2023selfrefine} 通过自我反馈迭代精炼输出；
STaR~\cite{zelikman2022star} 用模型自生成的推理路径构造监督信号；
Self-Improve~\cite{huang2023selfimprove} 以自一致性多数投票生成
高置信伪标注；Reflexion~\cite{shinn2023reflexion} 引入语言化自
我反思信号；Self-Rewarding~\cite{yuan2024selfrewarding} 与
SPIN~\cite{chen2024spin} 进一步将评判者与对手纳入自博弈闭环。这
些工作共同表明：一个足够强的模型可以充当自身的数据引擎。

将 RSI 范式迁移到推荐系统具有天然的吸引力：隐式反馈场景下"未观
测 $\neq$ 负样本"，模型对未观测交互的预测本身携带信息，自生成样
本相当于对模型知识的一种\emph{蒸馏回注}。对交互稀疏用户尤其如
此——冷启动与长尾场景正是监督信号最匮乏、自生成补充最可能产生
增益的位置（第四章的稀疏度分层分析将对这一直觉给出阴性检验结
果——假设的价值在于可证伪）。

\subsection{推荐系统中 RSI 的机会与风险}
从信息论视角看，RSI 闭环的增益潜力有明确上界：自生成样本携带的
新信息不超过模型已有知识（参数中编码的分布）与真实交互分布之间
的互补部分。这给出两个推论。其一，\emph{收益位置}：收益集中在
"模型已知但未被充分利用"的区域——交互稀疏的用户（其偏好未被充
分训练）、长尾物品（其嵌入欠拟合）、以及正则化不足的决策边界。
其二，\emph{风险位置}：风险集中在"模型错得最自信"的区域——头
部偏好的自我复制、以及训练分布尾部的人为抹平。机会与风险在
空间上高度重叠，这正是 RSI 控制问题的微妙之处：不能因风险而
关闭整个闭环（放弃增益），也不能因机会而放任准入（引发崩塌）
——需要的是对自生成数据的\emph{外科手术式}准入控制。

从系统演化视角看，RSI 闭环还是一类新的训练算法设计模式：传统
训练算法的性能曲线由固定数据集决定，而 RSI 系统的性能曲线由
"数据—模型"共演化动力学决定。这使得评估协议本身需要变化——
单点评估（最终指标）不足以刻画系统，必须评估\emph{轨迹}（指标
随迭代轮次的形态：收敛、震荡、发散）。本文第 4 章的实验设计
（迭代曲线 + 漂移曲线 + 跨轮稳定性指标）正是对这一评估哲学的
落实。

然而，RSI 闭环并非免费午餐。自我生成的训练信号存在系统性风险：
模型倾向于强化自己已经相信的预测，而头部物品在预测分布中天然占
优，于是"生成—训练—再生成"的闭环会形成\textbf{头部自强化正反
馈}，随迭代轮次放大为分布漂移与性能崩塌；若再叠加对生成样本质量
缺乏量化控制，低置信样本的持续混入将造成误差累积。如何在推荐场
景中\emph{控制}自生成数据的保真度、使 RSI 闭环安全可控而非引发
崩塌，是本文研究的核心问题。

\section{国内外研究现状概述}
与本工作相关的研究可粗分为四条线。（1）\textbf{序列推荐}：
GRU4Rec~\cite{hidasi2016gru4rec} 开创了会话级循环建模范式，
SASRec~\cite{kang2018sasrec} 与 BERT4Rec~\cite{sun2019bert4rec}
引入自注意力结构，成为当前序列推荐的主流骨干。（2）\textbf{自监
督与数据增强推荐}：SGL~\cite{xie2022sgl}、SimGCL~\cite%
{yu2022simgcl} 以对比学习挖掘自监督信号；
CoSeRec~\cite{liu2021coserec} 系统研究了序列扰动增强。
（3）\textbf{半监督学习中的自训练}：从
Yarowsky~\cite{yarowsky1995wsd} 的经典自训练到伪标签
~\cite{lee2013pseudo}、一致性正则
Mean-Teacher~\cite{tarvainen2017mean}、MixMatch~\cite%
{berthelot2019mixmatch} 与 FixMatch~\cite{sohn2020fixmatch}，置
信门控与增强一致性已被证明是控制伪标签质量的有效手段。
（4）\textbf{语言模型的递归自改进}：Self-Refine~\cite%
{madaan2023selfrefine}、STaR~\cite{zelikman2022star}、
Reflexion~\cite{shinn2023reflexion}、Self-Rewarding~\cite%
{yuan2024selfrewarding}、SPIN~\cite{chen2024spin} 构成 RSI 的方
法谱系。上述四条线的详细梳理与批判性分析见第二章。

值得注意的是，现有工作在"推荐系统中的 RSI 闭环保真度控制"上存
在明显空白：半监督学习的门控手段未与推荐场景的流行度偏差特性结
合，而语言模型的 RSI 方法未处理推荐所特有的分布漂移模式。本文正
是在这一空白处展开。

\section{研究问题}\label{sec:questions}
本文聚焦 RSI 闭环在序列推荐中的两个具体缺陷：

\begin{description}
\item[Q1 自生成数据分布漂移。] 模型自生成样本天然偏向其头部预
测，随迭代轮次训练分布持续偏离真实交互分布；这一漂移能否被量化
度量并在训练过程中被在线控制？
\item[Q2 迭代误差累积与崩塌。] 朴素自训练闭环为何崩塌、崩塌的动
力学是什么；何种生成与准入机制能将崩塌转化为稳定甚至正向的迭代
轨迹？
\end{description}

\section{研究内容与创新点}
针对上述问题，本文提出\textbf{保真度门控的递归自改进序列推荐框
架}，创新点如下：

\begin{enumerate}
\item \textbf{RSI 闭环的形式化与失败模式复现}（第
\ref{ch:framework} 章）：将 RSI 闭环形式化为五段迭代过程，给出分
布漂移与误差累积的度量；在 MovieLens-100K 上复现朴素自训练的崩
塌（NDCG@10 由 \ExpNaiveStart{} 崩塌至 \ExpNaiveEnd{}），并实证
其头部自强化根源。
\item \textbf{分层校准的保真度门控机制}：置信轴按流行度分层对预
测裕度做 z 标准化（实测保护的主要来源之一）；分布轴以分层 JS
距离预算约束混合分布漂移，作为候选池退化时才接管的条件性安全
网（第 4 章机制网格如实报告其在本数据集上从未咬合、净效应为
零，其约束力由单元级构造测试验证）。
\item \textbf{多视图一致性伪标签生成}：以 $K$ 个扰动视图上预测的
一致性为生成条件，使自生成样本编码扰动鲁棒信念而非头部偏好，为
闭环提供非自强化的信息来源。
\item \textbf{可复现实验基础设施}：全部实验以固定种子、单线程确
定性的方式实现，并以指令流编排（生成—门控—混合—重训—评估的
每轮各为一条可审计指令），实验结果可由命令行逐字节复现。
\end{enumerate}

\section{论文组织结构}
全文共五章。第一章绪论，给出背景、问题与创新点；第二章相关基础
与批判性综述；第三章问题形式化与框架设计，含机制推导与复杂度分
析；第四章实验设计与结果分析，含主对比、消融、迭代过程分析与分
布漂移可视化；第五章总结与展望。
